%==================================================================================================================
%================================================ Perspectives ====================================================
%==================================================================================================================

% --------------------------------- Perspectives : Complémenter et construire -------------------------------------
\section{Perspectives}
\begin{frame}[t]{Perspectives : compléter et construire}
  \vspace{0.4em}

  \begin{columns}[T,onlytextwidth]
    \column{0.48\textwidth}
      \begin{pcbanner}{Brique photovoltaïque}
        \begin{itemize}\footnotesize
          \item Ajouter les hétérojonctions
          \item Revoir les conditions aux limites
          \item Compléter la partie optique 
          \item Valider sur les données de l'ANR
          \item Intégrer des phénomènes fins si nécessaire
        \end{itemize}
      \end{pcbanner}
    \column{0.48\textwidth}
      \begin{pcbanner}{Brique Electrochimique}
        \begin{itemize}\footnotesize
          \item Partie hors du milieu poreux 
          \item Modèle GDE
          \item Vérifier la loi de vitesse (données ANR)
          \item Valider sur les données de l'ANR
          \item Intégrer des phénomènes fins si nécessaire (HER, pH, double couche\dots)
        \end{itemize}
      \end{pcbanner}
  \end{columns}

  \vspace{0.35em}
  \begin{columns}[T,onlytextwidth]
    \column{\textwidth}
      \begin{pcbanner}{Construire/Comprendre de nouvelles briques}
        \begin{itemize}\footnotesize
          \item Phototgénération : passer à l'electromagnétique
          \item \textit{Detailed balance} : approche alternative 
        \end{itemize}
      \end{pcbanner}
  \end{columns}
\end{frame}

% ---------------------------------- Perspectives : Coupler et optimiser -------------------------------------
\begin{frame}[t]{Perspectives : coupler et optimiser}
  \vspace{0.4em}

  \begin{pcbanner}{Coupler les briques}
    \begin{itemize}\footnotesize
      \item Articuler les briques une fois complétées et validées (notamment les CL)
      \item Commencer par l'interface semi-conducteur / électrolyte
    \end{itemize}
  \end{pcbanner}

  \vspace{0.8em}

  \begin{pcbanner}{Objectif final : Outil d'analyse d'extrapolation et d'optimisation}
    \begin{itemize}\footnotesize
      \item Modèle couplé et prédictif $\rightarrow$ explorer l'espace des géométries
      \item Épaisseurs, structure poreuse, aire catalytique
      \item Maximiser le courant et le rendement de production de \ch{CO}
    \end{itemize}
  \end{pcbanner}
\end{frame}


%==================================================================================================================
%================================================ Fin Perspectives ================================================
%==================================================================================================================
